\subsection{Modules over a product of rings}

Let \(\mathbf{A}\) be a ring and \((\mathfrak{b}_i)_{i\in I}\) a direct decomposition of \(\mathbf{A}\) , that is (I, p.~110, Def. 7) a finite family of two-sided ideals of \(\mathbf{A}\) such that the natural homomorphism from \(\mathbf{A}\) into the product of the \(\mathbf{A} / \mathfrak{b}_i\) is bijective. By loc. cit., Prop. 10, there exists a family \((e_i)_{i\in I}\) of central idempotents of \(\mathbf{A}\) such that \(\mathfrak{b}_i = \mathbf{A}(1 - e_i)\) , \(\sum_{i\in I}e_i = 1\) and \(e_i e_j = 0\) for \(i\neq j\) .

For every left A-module M, let \(M_{i}\) denote the set of \(m \in M\) such that \(b_{i}m = 0\) ; since \(b_{i}\) is a two-sided ideal, this is a submodule of M; moreover, if \(a, b \in A\) and \(a - b \in b_{i}\) , then the homotheties \(a_{M_{i}}\) and \(b_{M_{i}}\) coincide; there is thus a unique \((\mathbf{A}/\mathbf{b}_{i})\) -module structure on \(M_{i}\) such that the A-module structure of \(M_{i}\) is induced via the homomorphism \(A \to A/b_{i}\) .

PROPOSITION 1. --- The A-module M is the direct sum of its submodules \(M_{i}\) .

Let \(p_i: \mathbf{M} \to \mathbf{M}\) denote the homothety \(m \mapsto e_i m\) ; the map \(p_i\) is A-linear since \(e_i\) is central; since \(e_i^2 = e_i\) , \(\sum_{i \in I} e_i = 1\) and \(e_i e_j = 0\) for \(i \neq j\) , we have

\[
p _ {i} \circ p _ {i} = p _ {i}, \quad \sum_ {i \in I} p _ {i} = 1 _ {\mathrm{M}}, \quad p _ {i} \circ p _ {j} = 0 \quad \text { for } \quad i \neq j,
\]

and the \(p_i\) form an orthogonal family of projectors whose sum is the identity (II, p.~209, Def. 7). By loc. cit., Prop. 12, the module \(\mathbf{M}\) is the direct sum of the submodules \(p_i(\mathbf{M}) = e_i\mathbf{M}\) . In addition \(e_i\mathbf{M}\) is annihilated by \(\mathfrak{b}_i = \mathbf{A}(1 - e_i)\) ; if \(i \neq j\) and \(m \in \mathbf{M}\) then \((1 - e_i)e_jm = e_jm\) , whence no nonzero element of \(e_j\mathbf{M}\) is annihilated by \(1 - e_i\) , and so a fortiori by \(\mathfrak{b}_i\) . It follows that \(e_i\mathbf{M} = \mathbf{M}_i\) , and the proposition is proved.

Remarks. --- 1) Conversely, let \(\mathbf{M}_i'\) be an \((\mathbf{A} / \mathfrak{b}_i)\) -module for each \(i\) , and consider the A-module \(\mathbf{M}\) , the direct sum of the A-modules \(\mathbf{M}_i'\) ; then the submodules \(\mathbf{M}_i\) constructed above coincide with the \(\mathbf{M}_i'\) (it is enough to note that if \(i \neq j\) then no nonzero element of \(\mathbf{M}_j'\) is annihilated by \(\mathfrak{b}_i\) because \(\mathfrak{b}_i + \mathfrak{b}_j = \mathbf{A}\) ). Thus, roughly speaking, it amounts to the same thing to consider an A-module \(\mathbf{M}\) or a family \((\mathbf{M}_i)\) of modules over the rings \(\mathbf{A} / \mathfrak{b}_i = \mathbf{A}_i\) .

\begin{enumerate}
\def\labelenumi{\arabic{enumi})}
\setcounter{enumi}{1}
\item
  By the above proof, the projectors of M onto the components \(M_{i}\) are homotheties.
\item
  The A-module \(\mathbf{M}\) is cyclic if and only if each \(\mathbf{M}_i\) is cyclic: if \(\mathbf{M} = \mathbf{A}m\) , then \(\mathbf{M}_i = \mathbf{A}_i e_i m\) ; conversely if \(\mathbf{M}_i = \mathbf{A}_i m_i\) and \(m = \sum_{i \in I} m_i\) , then \(\mathbf{M} = \mathbf{A}m\) ; indeed, if
\end{enumerate}

\(n \in M\) projects onto \(a_i m_i\) for each \(i\) , and if \(a \in A\) is congruent to \(a_i\) mod \(b_i\) for each \(i\) , then \(am\) and \(n\) have the same image in each \(M_i\) , so coincide.

\begin{enumerate}
\def\labelenumi{\arabic{enumi})}
\setcounter{enumi}{3}
\tightlist
\item
  Let \(\mathbf{M}\) and \(\mathbf{N}\) be two A-modules, with components \((\mathbf{M}_i)\) and \((\mathbf{N}_i)\) . Let \(u \in \operatorname{Hom}_{\mathbf{A}}(\mathbf{M}, \mathbf{N})\) be an A-linear map from \(\mathbf{M}\) to \(\mathbf{N}\) ; then for all \(i\) and for all \(m \in \mathbf{M}_i\) we have \(u(m) \in \mathbf{N}_i\) , so \(u\) induces an \(\mathbf{A}_i\) -linear map \(u_i \in \operatorname{Hom}_{\mathbf{A}_i}(\mathbf{M}_i, \mathbf{N}_i)\) . It is easy to check that the map \(u \mapsto (u_i)\) is an isomorphism of \(\mathbf{Z}\) -modules (resp. of A-modules when \(\mathbf{A}\) is commutative)
\end{enumerate}

\[
\operatorname{Hom} _ {\mathbf {A}} (\mathbf {M}, \mathbf {N}) \rightarrow \prod_ {i \in I} \operatorname{Hom} _ {\mathbf {A} _ {i}} (\mathbf {M} _ {i}, \mathbf {N} _ {i}).
\]

